\section{Evidence-Preserving Head Steering}
\label{sec:method}

\begin{figure*}[t]
\centering
\begin{tikzpicture}[
  box/.style={draw=cvprblue,rounded corners=2pt,fill=cvprblue!5,
    align=center,text width=0.28\textwidth,minimum height=2.1cm,
    inner sep=6pt,font=\small},
  link/.style={-{Latex[length=2.5mm]},thick,color=cvprblue}]
\node[box] (discover) {
  \textbf{1. Localize heads}\\[3pt]
  Paired development images\\
  Target patch effect\\
  minus background-control effect};
\node[box,right=0.025\textwidth of discover] (fit) {
  \textbf{2. Fit bounded gains}\\[3pt]
  Separate balanced Yes/No data\\
  Frozen backbone; few scalars\\
  Preserve correct-answer distribution};
\node[box,right=0.025\textwidth of fit] (infer) {
  \textbf{3. Infer from one image}\\[3pt]
  Fixed head indices and gains\\
  Normal cached generation\\
  No donor, masks, or test labels};
\draw[link] (discover) -- (fit);
\draw[link] (fit) -- (infer);
\end{tikzpicture}
\caption{The implemented \method pipeline. Localization and fitting use
disjoint development images. Patching is an offline oracle diagnostic,
whereas deployment applies learned scalar gains to a single input.
Semantic-context robustness is an evaluation hypothesis, not a property
guaranteed by the localization score.}
\label{fig:pipeline}
\end{figure*}

\subsection{Task and intervention site}
Given an image $x$ and a question $q$ about the existence of a target
category, let $m(x,q)=z_{\mathrm{Yes}}-z_{\mathrm{No}}$ denote the base model's
first-answer-token logit margin. The configured affirmative and negative
strings must each correspond to one distinct token. This forced-choice
quantity is used for localization and fitting; it is distinct from the
accuracy of freely generated answers.

At decoder layer $\ell$, let $u_{\ell,h,t}$ be the output of query head $h$
at position $t$, before the attention output projection mixes heads. The
implementation hooks the concatenated input of that projection and selects
the corresponding head slice. For grouped-query attention, the split uses
the number of query heads, not the number of key/value heads. We do not
modify attention probabilities or claim that the selected head is an
independent semantic feature.

\subsection{Control-adjusted causal localization}
\label{sec:localize}
Each discovery record provides a clean image $x$, a target-occluded image
$x^{-}$, a background-control image $x^{c}$, and the same question $q$.
The current COCO builder unions the segmentation masks of all annotated
instances of the target category and fills the union with gray. A control
is created by translating the same mask shape onto background that excludes
annotated objects. Records without a valid matched control are unsuitable
for this analysis. Annotation incompleteness and the limits of occlusion
as an intervention must be assessed separately.

We retain positive records for which the base model's clean forced-choice
answer is correct, $m(x,q)>0$, and target occlusion decreases the margin by
$d=m(x,q)-m(x^{-},q)>\epsilon$. For each candidate head, we replace its final
prompt-position activation in $x^{-}$ with the corresponding activation
from $x$, and repeat this operation for $x^{c}$. Writing the resulting
margin changes as $r_{\ell,h}$ and $c_{\ell,h}$, respectively, the score is
\begin{equation}
 s_{\ell,h}(x,q)=
 \frac{r_{\ell,h}-\lambda_c |c_{\ell,h}|}
      {\max(d,\eta)}.
 \label{eq:score}
\end{equation}
The denominator floor $\eta$ avoids treating tiny margin changes as strong
evidence. The control penalty discourages heads whose clean substitution
also strongly changes the matched control input. Scores are averaged over
questions from each original image, then aggregated by their median across
image groups. The top $K$ heads form the selected set $\mathcal{S}$.

The current selector always returns the top $K$ candidates; it does not
require positive scores. Therefore, selection alone must not be interpreted
as proof that useful heads were discovered. We inspect score signs,
between-image variability, and stability across discovery samples before
interpreting a selected set. Random-head and clean-activation-RMS selectors
use the same candidate pool for controlled comparisons.

\paragraph{What this intervention establishes.}
A clean donor carries information about an object that the occluded input
does not contain. Successful patching can consequently increase an
\emph{incorrect} positive answer on the occluded image. Its role is to
identify a route by which the clean input affects the answer, not to count
a corrected prediction. The background control only adjusts for some
occlusion-related effects. It neither identifies a unique causal pathway
nor proves invariance to semantic context. The usefulness of
Eq.~\eqref{eq:score} is tested by subsequent single-input intervention,
not assumed.

\subsection{Bounded gains with distribution preservation}
\label{sec:fit}
On a separate fitting split, we freeze the pretrained model and optimize
one scalar $a_{\ell,h}$ per selected head. Its bounded gain is
$\delta_{\ell,h}=\rho\tanh(a_{\ell,h})$, initialized at zero, and its output
is transformed as
\begin{equation}
 \widetilde u_{\ell,h,t}=
 \begin{cases}
 (1+\delta_{\ell,h})u_{\ell,h,t},
       &(\ell,h)\in\mathcal{S},\\
 u_{\ell,h,t},&\text{otherwise}.
 \end{cases}
 \label{eq:gain}
\end{equation}
With $\rho=0.25$, the multiplier remains between $0.75$ and $1.25$.
The intervention cannot reverse or eliminate a head. Both amplification
and attenuation are learned; a high patch score does not prescribe their
sign.

Let $y\in\{-1,+1\}$ be the object-existence label, $m_\delta$ the steered
margin, and $p_0$ and $p_\delta$ the full-vocabulary next-token distributions
at the final prompt position. The fitting objective is
\begin{equation}
\begin{split}
\mathcal{L}={}&
 \mathbb{E}_{\mathrm{bal}}\left[
 \operatorname{softplus}(-y m_\delta)
 +\lambda_{\mathrm{KL}} I_{\mathrm{keep}}
   D_{\mathrm{KL}}(p_0\Vert p_\delta)\right]\\
 &+\frac{\lambda_2}{K}
   \sum_{(\ell,h)\in\mathcal{S}}\delta_{\ell,h}^{\,2}.
\end{split}
\label{eq:loss}
\end{equation}
Balanced sampling pairs positive and negative examples, cycling the smaller
class to give both labels equal exposure. $I_{\mathrm{keep}}$ is one only
when the base forced-choice answer is correct and the record is eligible
for preservation. Eligibility is fixed by the data protocol, not chosen
after inspecting intervention outcomes.

The KL term regularizes a single next-token distribution; it does not
guarantee preservation of a generated caption or general reasoning.
Those properties require held-out evaluation. Although the backbone is
frozen, fitting uses supervision and backpropagation, so \method is not
training-free.

\subsection{Deployment and computational requirements}
The default intervention affects only the final prefill token and the
current token at each cached decoding step. Earlier prompt positions are
unchanged. An all-prefill-position variant is a separate ablation, and its
setting must match between fitting and deployment. The deployed artifact
contains the head indices, gains, and intervention-position setting.

Inference uses the original image and question, without edited images,
clean donors, ground-truth labels, or a per-example backward pass. It uses
one normal generation path, but hooks and head-wise operations still incur
overhead; no latency benefit is assumed. Localization is more expensive:
each discovery record requires three base forwards and each retained
record requires two additional forwards per candidate head. Fitting also
requires backward computation through the frozen network. We report these
offline costs separately from inference cost.

\subsection{Semantic-context robustness as a testable hypothesis}
The implemented method does not optimize an explicit semantic-context
invariance loss. Its proposed connection to context robustness is that
target-selective localization may expose intervention sites that help
preserve useful visual evidence. This connection is falsifiable. If gains
selected by Eq.~\eqref{eq:score} fail on semantic-context edits, or perform
no better than gains on arbitrary heads, the motivating mechanism is
unsupported even if a standard hallucination score improves.
